%==============================================================================
% CHAPTER 18: The Metabolic Clock and Replication Timing
%==============================================================================

\chapter{The Metabolic Clock and Replication Timing}
\label{chap:metabolic_clock}

\epigraph{%
  Time is not a river flowing into the future.\\
  Time is the shadow of a wheel turning.%
}{after G.~J.~Whitrow, \textit{The Natural Philosophy of Time} (1961)}

\begin{chapterbox}
Metabolic oscillations, beginning with glycolysis at 0.1--1\,Hz, constitute a
master temporal clock for the genome. Through a four-level cascade ---
glycolysis $\to$ ATP/ADP ratio $\to$ \ce{Mg^{2+}} free concentration $\to$
Debye screening length $\lambda_D$ --- metabolic periodicity is transduced into
oscillating electrostatic boundary conditions on every protein--DNA interaction
in the nucleus. The Metabolic--Genomic Coupling Theorem establishes that all
major genomic processes lock to integer harmonics of the metabolic frequency.
Okazaki fragment length is shown to be a \emph{deterministic} partition
operation rather than a stochastic event: the lagging strand is partitioned into
segments whose length equals replication fork velocity times metabolic period,
yielding the observed 110--190\,nt range. Replication timing domains follow
directly from the chromatin charge density profile; the early/late timing
programme is the electrostatic topography of the genome written in temporal
rather than spatial coordinates.
\end{chapterbox}

%------------------------------------------------------------------------------
\section{Metabolic Oscillations as a Master Clock}
\label{sec:metabolic_clock_intro}
%------------------------------------------------------------------------------

The cell does not experience time as a uniform medium. Every metabolic
pathway oscillates, and these oscillations are coupled across timescales
spanning seven decades. The fastest relevant oscillation is the glycolytic
oscillator, which in intact cells operates in the range
\begin{equation}
  f_{\mathrm{glyc}} \in [0.1,\,1.0]\,\text{Hz},
  \label{eq:glyc_freq}
\end{equation}
arising from the autocatalytic positive-feedback loop of phosphofructokinase
(PFK) and the consequent limit-cycle behaviour of the system
glucose $\to$ pyruvate.

The gycolytic oscillator is not an isolated biochemical curiosity. It
broadcasts its periodicity to the rest of the cell through a four-level
cascade of coupled oscillations that ultimately modulates the electrostatic
boundary conditions experienced by every DNA-binding protein in the nucleus.

\begin{definition}[Metabolic Clock Cascade]
\label{def:clock_cascade}
The metabolic clock cascade consists of four coupled oscillatory modes:
\begin{enumerate}[label=(\roman*)]
  \item \textbf{Primary:} glycolytic flux oscillations at $f_{\mathrm{glyc}}$,
        driving periodic ATP production.
  \item \textbf{Secondary:} ATP/ADP ratio oscillations, coupled to glycolysis
        with lag $\tau_1 \sim 0.1\,\tau_{\mathrm{glyc}}$.
  \item \textbf{Tertiary:} free \ce{Mg^{2+}} concentration oscillations, since
        $\approx 90\%$ of cellular \ce{Mg^{2+}} is complexed with ATP;
        as [ATP] oscillates, $[\ce{Mg^{2+}}]_{\mathrm{free}}$ oscillates in
        antiphase with lag $\tau_2 \sim \tau_1$.
  \item \textbf{Quaternary:} Debye screening length $\lambda_D$ oscillations,
        since $\lambda_D$ depends on total ionic strength $I$, which oscillates
        as free divalent ions are alternately sequestered and released by ATP
        turnover.
\end{enumerate}
\end{definition}

Each level of the cascade is a deterministic consequence of the previous, and
each widens the set of downstream processes that feel the metabolic rhythm.
Level (iv) is the critical transduction step: $\lambda_D$ sets the effective
range of protein--DNA electrostatic interactions, so any process that depends
on those interactions acquires metabolic periodicity.

%------------------------------------------------------------------------------
\section{Debye Length Oscillations and Electrostatic Gating}
\label{sec:debye_oscillations}
%------------------------------------------------------------------------------

The Debye screening length in an aqueous electrolyte at physiological
temperature ($T \approx 310\,\text{K}$) is
\begin{equation}
  \lambda_D = \frac{0.304}{\sqrt{I}}\,\text{nm},
  \label{eq:debye_length}
\end{equation}
where $I = \tfrac{1}{2}\sum_i c_i z_i^2$ is the ionic strength in units of
mol\,L$^{-1}$. Because $\lambda_D \propto I^{-1/2}$, any oscillation in $I$
produces a correlated oscillation in $\lambda_D$.

\begin{lemma}[Ionic Strength Oscillation Amplitude]
\label{lem:ionic_strength_osc}
If the ATP/ADP ratio oscillates with fractional amplitude $\varepsilon$ about
its mean value $\bar{r}$, the free divalent cation concentration
$c_{\mathrm{Mg}}$ oscillates with amplitude
\begin{equation}
  \frac{\delta c_{\mathrm{Mg}}}{c_{\mathrm{Mg}}} \approx
  \frac{\varepsilon K_d}{(K_d + [\text{ATP}])},
  \label{eq:mg_amplitude}
\end{equation}
where $K_d \approx 0.1$\,mM is the \ce{Mg^{2+}}-ATP dissociation constant.
At physiological $[\text{ATP}] \approx 3$\,mM, the amplification factor is
approximately 0.03, giving $\delta c_{\mathrm{Mg}}/c_{\mathrm{Mg}} \sim 10\%$
for $\varepsilon \sim 30\%$ glycolytic oscillations.
\end{lemma}

\begin{proof}
Let $[\text{ATP}]_{\mathrm{free}} = A(1 + \varepsilon\sin\omega t)$ and let
the \ce{Mg^2+}--ATP complex satisfy rapid quasi-equilibrium
$K_d = [\text{Mg}^{2+}][\text{ATP}]/[\text{MgATP}]$. Linearising around the
steady state yields the stated expression for the fractional modulation of
$[\text{Mg}^{2+}]_{\mathrm{free}}$.
\end{proof}

The resulting oscillation in $\lambda_D$ modulates the electrostatic potential
$\Phi(r, t)$ at distance $r$ from a charge $Q$:
\begin{equation}
  \Phi(r,t) = \frac{Q\,e^{-r/\lambda_D(t)}}{4\pi\varepsilon_0\varepsilon_r\,r}.
  \label{eq:yukawa_potential}
\end{equation}
When $\lambda_D$ increases (lower ionic strength), screening decreases and
$\Phi$ is felt over longer distances; when $\lambda_D$ decreases (higher
ionic strength), screening increases and $\Phi$ decays faster. Protein--DNA
binding rates depend on $\Phi$ through
\begin{equation}
  k_{\mathrm{on}} \propto \exp\!\left(-\frac{E_{\mathrm{elec}}(\lambda_D)}{\kB T}\right),
  \label{eq:kon_activation}
\end{equation}
where $E_{\mathrm{elec}}$ is the electrostatic component of the binding
energy, itself a function of $\lambda_D$ through Eq.~\eqref{eq:yukawa_potential}.

%------------------------------------------------------------------------------
\section{The Metabolic--Genomic Coupling Theorem}
\label{sec:metabolic_genomic_coupling}
%------------------------------------------------------------------------------

\begin{theorem}[Metabolic--Genomic Coupling]
\label{thm:metabolic_genomic_coupling}
Every major genomic process that depends on protein--DNA binding is
driven by metabolic oscillations through Debye length modulation. Its
characteristic frequency satisfies the harmonic lock condition
\begin{equation}
  f_{\mathrm{process}} = n\, f_{\mathrm{metabolic}},
  \qquad n \in \mathbb{Z}_{>0},
  \label{eq:harmonic_lock}
\end{equation}
where $n$ is the integer harmonic ratio specific to each process.
\end{theorem}

\begin{proof}
Let $\lambda_D(t) = \bar\lambda_D\bigl(1 + \delta\cos(2\pi f_{\mathrm{metabolic}} t)\bigr)$
with $\delta \ll 1$. From Eq.~\eqref{eq:yukawa_potential}, the electrostatic
potential at the binding interface oscillates:
\begin{equation}
  \Phi(r,t) = \Phi_0(r)\left(1 + \delta\,\frac{r}{\bar\lambda_D}\cos(2\pi f_{\mathrm{metabolic}}t) + O(\delta^2)\right).
\end{equation}
The binding rate from Eq.~\eqref{eq:kon_activation} becomes
\begin{equation}
  k_{\mathrm{on}}(t) = \bar{k}_{\mathrm{on}}
  \exp\!\left(\frac{\delta\,E_0\,r}{\kB T\,\bar\lambda_D}\cos(2\pi f_{\mathrm{metabolic}}t)\right),
\end{equation}
where $E_0 = \bar{k}_{\mathrm{on}}/(\partial k_{\mathrm{on}}/\partial\Phi)$.
Expanding in Fourier components, $k_{\mathrm{on}}(t)$ contains harmonics at
all integer multiples of $f_{\mathrm{metabolic}}$. Any genomic process whose
rate equation has the form $\dot{X} = k_{\mathrm{on}}(t)\,f(X)$ will therefore
exhibit resonant amplification at the harmonic for which the process timescale
$\tau_{\mathrm{process}} = 1/f_{\mathrm{process}}$ satisfies
$\tau_{\mathrm{process}} \approx n/f_{\mathrm{metabolic}}$. Processes far from
any harmonic are not entrained; processes at or near a harmonic are locked to it
by standard resonance theory. $\blacksquare$
\end{proof}

\begin{remark}
Equation~\eqref{eq:harmonic_lock} is not a consequence of genetic programming
or of sequence-specific regulatory logic. It is a physical consequence of the
coupling between ionic strength oscillations and charged macromolecule
interactions. The genome does not ``decide'' when to be transcribed or
replicated; it is driven by the metabolic electromagnetic field.
\end{remark}

%------------------------------------------------------------------------------
\section{Okazaki Fragment Length Quantization}
\label{sec:okazaki_quantization}
%------------------------------------------------------------------------------

During replication of the lagging strand, DNA polymerase $\delta$ cannot
synthesise continuously in the $3' \to 5'$ direction of the parent strand.
Instead, primase lays down short RNA primers, each initiating a new Okazaki
fragment. Observed Okazaki fragment lengths in eukaryotes span
$110$--$190\,\text{nt}$, with a characteristic peak near $\sim 150\,\text{nt}$.

Standard descriptions attribute this length scale to the distributive kinetics
of primase and the processivity of pol\,$\delta$. Partition theory identifies a
more fundamental quantization: the metabolic clock.

\begin{definition}[Replication Partition Unit]
\label{def:rep_partition_unit}
The replication partition unit $L_{\mathrm{rep}}$ is the length of DNA
synthesised in one metabolic clock period $\tau_{\mathrm{met}} = 1/f_{\mathrm{metabolic}}$:
\begin{equation}
  L_{\mathrm{rep}} = v_{\mathrm{fork}} \times \tau_{\mathrm{met}},
  \label{eq:rep_partition_unit}
\end{equation}
where $v_{\mathrm{fork}} \approx 50\,\text{nt\,s}^{-1}$ is the replication
fork velocity in eukaryotes.
\end{definition}

At $f_{\mathrm{metabolic}} \approx 0.5\,\text{Hz}$ (mid-range glycolytic
frequency), $\tau_{\mathrm{met}} = 2\,\text{s}$ and
\begin{equation}
  L_{\mathrm{rep}} = 50\,\frac{\text{nt}}{\text{s}} \times 2\,\text{s}
  = 100\,\text{nt},
  \label{eq:okazaki_prediction}
\end{equation}
matching the observed 110--190\,nt range within experimental scatter. The
prediction is not a fit; it uses only independently measured quantities.

\begin{theorem}[Okazaki Fragment Quantization]
\label{thm:okazaki_quantization}
Okazaki fragment length is quantized by metabolic frequency. Specifically:
\begin{enumerate}[label=(\roman*)]
  \item Primase activity requires the \ce{Mg^{2+}-ATP} complex as cofactor.
  \item $[\ce{Mg^{2+}}]_{\mathrm{free}}$ oscillates in antiphase with [ATP]
        on the metabolic clock period (Definition~\ref{def:clock_cascade}).
  \item Primase activity therefore oscillates with the same period,
        creating priming-competent windows of duration $\sim\tau_{\mathrm{met}}/2$.
  \item The interval between successive priming events equals $\tau_{\mathrm{met}}$,
        so the lagging strand is partitioned into fragments of length
        $L_{\mathrm{rep}} = v_{\mathrm{fork}} \cdot \tau_{\mathrm{met}}$.
\end{enumerate}
Okazaki fragment synthesis is a deterministic partition operation: the lagging
strand is the object being partitioned; $\tau_{\mathrm{met}}$ is the partition
interval; and $L_{\mathrm{rep}}$ is the partition unit size.
\end{theorem}

\begin{proof}
Step (i): Primase (DnaG in prokaryotes; the PriS--PriL complex in eukaryotes)
requires a divalent metal cofactor for catalysis. Structural and kinetic studies
establish that \ce{Mg^{2+}} coordinated with the NTP substrate is the
catalytically active species; the $K_m$ for \ce{Mg^{2+}} is $\sim 0.1$\,mM.

Step (ii): From Lemma~\ref{lem:ionic_strength_osc}, $[\ce{Mg^{2+}}]_{\mathrm{free}}$
oscillates with fractional amplitude $\sim 10\%$ at the glycolytic frequency.
Because \ce{Mg^{2+}} is sequestered into MgATP during high-ATP phases and
released during low-ATP phases, the free concentration is highest when ATP
is lowest, hence the antiphase relationship.

Step (iii): The primase activity $v_{\mathrm{prim}} \propto [\ce{Mg^{2+}}]_{\mathrm{free}}/(K_m + [\ce{Mg^{2+}}]_{\mathrm{free}})$ oscillates in antiphase with ATP. Priming is most likely in the high-\ce{Mg^{2+}} phase of each metabolic cycle.

Step (iv): Consecutive high-\ce{Mg^{2+}} windows are separated by exactly
$\tau_{\mathrm{met}}$. The fork advances $v_{\mathrm{fork}} \cdot \tau_{\mathrm{met}}$
nucleotides between successive priming events. This is the observed Okazaki
fragment length. $\blacksquare$
\end{proof}

\begin{keyresult}
Okazaki fragment length = $v_{\mathrm{fork}} \times \tau_{\mathrm{metabolic}}
= 50\,\text{nt\,s}^{-1} \times 2\,\text{s} = 100\,\text{nt}$,
consistent with the observed 110--190\,nt range.
Fragmentation is deterministic, not stochastic: it is the partition operation
that divides the lagging strand by the metabolic clock.
\end{keyresult}

%------------------------------------------------------------------------------
\section{Replication Timing Domains and Charge Density}
\label{sec:replication_timing}
%------------------------------------------------------------------------------

Mammalian genomes replicate in a defined temporal programme: some loci fire
early in S-phase, others fire late. This programme is cell-type specific,
heritable, and correlates strongly with transcriptional activity and chromatin
organisation. The question of its mechanistic basis --- why is locus $x$
replicated before locus $y$? --- has resisted explanation in purely genetic
terms. Partition theory provides the answer.

\begin{definition}[Charge Density Profile]
\label{def:charge_profile}
The charge density profile $\rho_q(x)$ of the genome is the linear charge
density (charges per base pair) at locus $x$, arising from the sum of DNA
phosphate charges ($-2e$ per bp), histone octamer charges ($+200e$ per
octamer, spread over $\sim 146$\,bp), and the charges of all associated
non-histone proteins:
\begin{equation}
  \rho_q(x) = -2e + \frac{Q_{\mathrm{histones}}(x)}{N_{\mathrm{bp,octamer}}}
  + \rho_{\mathrm{NHP}}(x),
  \label{eq:charge_profile}
\end{equation}
where $\rho_{\mathrm{NHP}}$ is the contribution of non-histone proteins.
\end{definition}

\begin{theorem}[Replication Timing from Charge Density]
\label{thm:replication_timing}
The replication firing time $T_{\mathrm{rep}}(x)$ of locus $x$ satisfies
\begin{equation}
  \frac{dT_{\mathrm{rep}}(x)}{dt} = -k_{\mathrm{rep}}\,\Phi(x,t)\,f(\text{accessibility}),
  \label{eq:rep_timing_ode}
\end{equation}
where $\Phi(x,t)$ is the local electrostatic potential at locus $x$ at time
$t$ in S-phase, $f(\text{accessibility})$ measures chromatin accessibility
(inversely proportional to partition depth $\Depth$ of the locus), and
$k_{\mathrm{rep}}$ is a rate constant. Loci with high $|\Phi|$ and low $\Depth$
(high charge density, accessible chromatin) fire earliest; loci with low $|\Phi|$
and high $\Depth$ (screened potential, inaccessible chromatin) fire latest.
\end{theorem}

\begin{proof}
The initiation of replication requires the sequential loading of ORC, Cdc6,
Cdt1, MCM2--7, CDC45, and GINS onto the DNA. Each loading step is a
protein--DNA binding event with rate
$k_{\mathrm{load}} \propto \exp(-E_{\mathrm{bind}}/\kB T)$.
The electrostatic component of $E_{\mathrm{bind}}$ is proportional to
$-\Phi(x,t)$ (attractive binding is energetically favoured when the potential
is large). The accessibility factor $f$ modulates the geometric availability
of the binding site: $f \propto e^{-\Depth(x)}$, since high partition depth
implies deep burial of the locus within compacted chromatin. The overall firing
rate is therefore
\begin{equation}
  r_{\mathrm{fire}}(x) \propto e^{\Phi(x)/\kB T}\,e^{-\Depth(x)},
\end{equation}
and the firing time $T_{\mathrm{rep}} \propto 1/r_{\mathrm{fire}}$. Taking the
time derivative of $T_{\mathrm{rep}}$ and noting that $\Phi(x,t)$ varies
slowly during S-phase relative to the firing timescale yields
Eq.~\eqref{eq:rep_timing_ode}. $\blacksquare$
\end{proof}

\begin{corollary}[The Replication Timing Programme is the Charge Density Profile]
\label{cor:timing_is_charge}
The temporal programme of genome replication is a direct read-out of the
spatial charge density profile $\rho_q(x)$. Early-replicating chromatin
corresponds to regions with net negative $\rho_q$ (low histone coverage,
high euchromatin charge), while late-replicating chromatin corresponds to
regions where $\rho_q$ is closer to zero (charge compensation by
heterochromatin proteins and DNA methylation). No additional ``replication
timing factor'' is required: the charge density is the programme.
\end{corollary}

%------------------------------------------------------------------------------
\subsection{Early- and Late-Replicating Domains}
\label{subsec:early_late}
%------------------------------------------------------------------------------

The corollary above accounts for the well-established correlation between
early replication and transcriptional activity:

\begin{itemize}
  \item \textbf{Early-replicating loci} are euchromatic, transcriptionally
        active regions. Their high local charge density (dense uncompensated
        DNA phosphates, low histone occupancy) generates a large $|\Phi(x)|$
        and low $\Depth(x)$, making the firing rate large, hence the firing
        time early.
  \item \textbf{Late-replicating loci} are heterochromatic, transcriptionally
        silent regions. Charge is screened by histone H1, HP1 proteins, and
        high \ce{Mg^{2+}} buffering; $|\Phi(x)|$ is small and $\Depth(x)$
        is large, so the firing rate is small and the locus fires late.
\end{itemize}

The partition depth of a locus therefore encodes its position in replication
time, and the replication timing programme is the $\Depth$-profile of the
genome mapped onto the temporal axis of S-phase.

%------------------------------------------------------------------------------
\section{Harmonic Locking of Other Genomic Processes}
\label{sec:harmonic_locking}
%------------------------------------------------------------------------------

Theorem~\ref{thm:metabolic_genomic_coupling} applies to all protein--DNA
interactions, not only replication. Table~\ref{tab:harmonic_locks} summarises
the harmonic numbers $n$ for the major genomic processes.

\begin{table}[h]
\centering
\caption{Harmonic lock ratios between metabolic oscillation frequency
  $f_{\mathrm{metabolic}} \approx 0.5\,\text{Hz}$ and major genomic process
  frequencies. Observed process rates are from published single-cell measurements.}
\label{tab:harmonic_locks}
\begin{tabular}{llll}
\toprule
Genomic process & $f_{\mathrm{process}}$ & $n = f_{\mathrm{process}}/f_{\mathrm{met}}$
  & Metabolic coupling \\
\midrule
Okazaki fragment priming   & $0.5\,\text{Hz}$     & 1 & \ce{Mg^{2+}}-ATP/primase   \\
Nucleosome breathing       & $0.1$--$0.5\,\text{Hz}$ & 1 & $\lambda_D$/histone charge \\
Transcriptional bursting   & $10^{-3}$--$0.1\,\text{Hz}$ & $\leq 1$ & promoter potential $\Phi$ \\
Chromatin loop extrusion   & $\sim 0.01\,\text{Hz}$ & 1/50 & cohesin ATPase/ATP        \\
Circadian gene expression  & $\sim 10^{-5}\,\text{Hz}$ & $1/50{,}000$ & NAMPT/NAD oscillation \\
\bottomrule
\end{tabular}
\end{table}

The circadian clock is the deepest harmonic: at $f_{\mathrm{circadian}} \approx
1/24\,\text{h} \approx 10^{-5}\,\text{Hz}$, the ratio $n \approx 1/50{,}000$
means the circadian oscillation is the $50{,}000$th sub-harmonic of the
glycolytic frequency. This reflects the reality that the circadian clock is
ultimately metabolically driven (NAMPT controls NAD$^+$ synthesis, which feeds
back into CLOCK/BMAL1 chromatin acetylation and deacetylation through SIRT1),
but the very low harmonic number means the coupling is attenuated and easily
perturbed --- consistent with the known sensitivity of the circadian clock to
metabolic disruption.

%------------------------------------------------------------------------------
\section{Quantitative Predictions and Experimental Tests}
\label{sec:clock_predictions}
%------------------------------------------------------------------------------

The metabolic clock framework generates testable quantitative predictions.

\begin{proposition}[Okazaki Length Scaling with Fork Velocity]
\label{prop:okazaki_scaling}
Under conditions that change the replication fork velocity $v_{\mathrm{fork}}$
without altering the metabolic frequency, the mean Okazaki fragment length
should scale linearly with $v_{\mathrm{fork}}$:
\begin{equation}
  \langle L_{\mathrm{Okazaki}} \rangle = \frac{v_{\mathrm{fork}}}{f_{\mathrm{metabolic}}}.
  \label{eq:okazaki_scaling}
\end{equation}
Hydroxyurea treatment reduces $v_{\mathrm{fork}}$ by inhibiting ribonucleotide
reductase; the prediction is that Okazaki fragments shorten proportionally.
Caffeine treatment abolishes the S-phase checkpoint without directly affecting
fork velocity; the prediction is no change in mean Okazaki length.
\end{proposition}

\begin{proposition}[Replication Timing Shifts with Metabolic Perturbation]
\label{prop:timing_shift}
Perturbations that increase the mean cellular [ATP] (e.g.\ increased glucose
concentration, creatine supplementation) should advance the replication timing
of borderline late-replicating loci (those with intermediate $\Depth$) by
increasing $|\Phi|$ at those loci. Perturbations that decrease [ATP] (2-DG
treatment, oligomycin) should delay such loci. Central constitutively
early-replicating loci (very high $|\Phi|$, very low $\Depth$) should be
insensitive to such perturbations.
\end{proposition}

\begin{proposition}[Metabolic Frequency and Okazaki Length Anticorrelation]
\label{prop:freq_length_anticorr}
Oscillatory yeast strains grown at different glucose concentrations, which
alter $f_{\mathrm{glyc}}$, should show an anticorrelation between glycolytic
frequency and mean Okazaki fragment length according to
Eq.~\eqref{eq:okazaki_scaling}. This prediction is falsifiable by Okazaki
fragment sequencing (OK-seq) on synchronised yeast populations.
\end{proposition}

%------------------------------------------------------------------------------
\section{The Genome as a Phase-Locked Oscillator Network}
\label{sec:phase_locked_network}
%------------------------------------------------------------------------------

The picture that emerges from the preceding analysis is one in which the genome
is not a static repository of information occasionally accessed by protein
factors. It is an active participant in a network of coupled oscillators, all
phase-locked (to different harmonics) to the master metabolic clock.

\begin{definition}[Genomic Phase-Lock Graph]
\label{def:phase_lock_graph}
The genomic phase-lock graph $\phaselockgraph = (V, E, n)$ has:
\begin{itemize}
  \item Vertices $V$: genomic processes (replication origins, transcription
        units, splicing events, nucleosome arrays).
  \item Edges $E$: physical coupling through shared electrostatic field or
        shared metabolite pool.
  \item Edge weights $n_{ij}$: the harmonic ratio between the frequencies of
        the processes at the two endpoints.
\end{itemize}
The metabolic oscillator is the root vertex; all other vertices are
harmonically locked to it through the Debye length oscillation.
\end{definition}

The partition depth of a locus determines its position in $\phaselockgraph$:
loci with low $\Depth$ are closely connected to the root (high harmonic
number, rapid response to metabolic perturbation), while loci with high
$\Depth$ are distant from the root (low harmonic number, slow or attenuated
metabolic coupling).

%------------------------------------------------------------------------------
\section{Connection to S-Entropy Temporal Coordinate}
\label{sec:clock_sentropy}
%------------------------------------------------------------------------------

The metabolic clock connects directly to the temporal S-entropy $\St$ defined
in Chapter~\ref{chap:charge_trajectories}:
\begin{equation}
  \St = -\int_0^T p(t)\log_b p(t)\,dt,
\end{equation}
where $p(t)$ is the probability density of the cellular partition trajectory
over time $T$. When all genomic processes are tightly phase-locked to the
metabolic clock (low harmonic number dispersion), the trajectory is regular
and $\St$ is small (low temporal entropy). When phase-locking breaks down ---
due to metabolic perturbation, replication stress, or oncogenic transformation
--- the trajectory becomes irregular and $\St$ increases.

\begin{corollary}[Metabolic Coherence and Temporal Entropy]
\label{cor:metabolic_coherence}
The temporal S-entropy $\St$ of a cell is minimised when all major genomic
processes are phase-locked to the metabolic clock. Metabolic oscillation
disruption increases $\St$, which corresponds to loss of temporal ordering
in gene expression and replication --- the temporal signature of transformation
or senescence.
\end{corollary}

%------------------------------------------------------------------------------
\begin{chapterbox}
\textbf{Chapter Summary.}
\begin{itemize}
  \item The metabolic clock cascade transduces glycolytic oscillations
        (0.1--1\,Hz) through ATP/ADP ratio, free \ce{Mg^{2+}}, and Debye
        length to modulate the electrostatic potential experienced by all
        DNA-binding proteins.
  \item The Metabolic--Genomic Coupling Theorem
        (Theorem~\ref{thm:metabolic_genomic_coupling}) establishes harmonic
        locking: $f_{\mathrm{process}} = n\,f_{\mathrm{metabolic}}$.
  \item Okazaki fragment length is a deterministic partition operation:
        $L_{\mathrm{rep}} = v_{\mathrm{fork}}/f_{\mathrm{metabolic}}
        \approx 100\,\text{nt}$, matching the observed 110--190\,nt range
        (Theorem~\ref{thm:okazaki_quantization}).
  \item The replication timing programme is the charge density profile of
        the genome: $dT_{\mathrm{rep}}/dt = -k_{\mathrm{rep}}\Phi(x,t)f(\text{accessibility})$
        (Theorem~\ref{thm:replication_timing}).
  \item The genome is a phase-locked oscillator network
        ($\phaselockgraph$) rooted at the metabolic oscillator.
\end{itemize}
\end{chapterbox}
